\subsection{Projective ideal evaluation and the intrinsic calculus}
\label{subsec:projective-ideal-bridge}

The normalized graded seminorm model of projective space carries a polynomial
Gauss extension of radius $\exp(-1)$. The following statements identify its
ideal functions with the intrinsic ideal calculus on affine seminorm spectra.
They concern the existing normalized model; they do not identify it with a
canonical projective analytification or construct coherent ideal sheaf gluing.

\begin{lemma}[Bounded Gauss extension {\cite[\S2.1]{FHX}}]
\label{projective:bounded-gauss-extension}
\lean{EnvelopeProperty.ProjBerk.gaussExt_isBounded}
\leanok
\uses{def:bounded-ideal-norm,def:pn-gauss}
For a normalized graded seminorm $v$ on
$R=k[x_0,\ldots,x_n]$, its Gauss extension $\sigma(v)$ is bounded on
$R[t]$: $|F|_{\sigma(v)}\le1$ for every $F\in R[t]$.
\end{lemma}
\begin{proof}
Every coordinate has value at most one and the coefficient field is trivially
valued, so the ultrametric polynomial estimate gives $|a|_v\le1$ for all
$a\in R$. The Gauss formula gives
$|F|_{\sigma(v)}=\max_j |F_j|_v\exp(-j)\le1$.
This is precisely boundedness on the coefficient ring $R[t]$.
\end{proof}

\begin{theorem}[Correspondence of projective ideal functions {\cite[\S2.1]{FHX}}]
\label{projective:ideal-spectrum-bridge}
\lean{EnvelopeProperty.ProjBerk.phiNorm_eq_idealNorm,
EnvelopeProperty.ProjBerk.phiIdeal_eq_idealLog,
EnvelopeProperty.ProjBerk.phiNorm_eq_spectrumIdealNorm,
EnvelopeProperty.ProjBerk.phiIdeal_eq_spectrumIdealLog}
\leanok
\uses{projective:bounded-gauss-extension,def:bounded-ideal-norm,
def:general-ideal-norm,spectrum:trivial-field-bridge}
For every ideal $I\subseteq R[t]$ and every normalized model point $v$,
\[
 \phi_I^{\mathrm{mult}}(v)=|I|_{\sigma(v)},\qquad
 \phi_I(v)=\log|I|(\sigma(v)).
\]
These equalities are unchanged when the prescribed trivial base values are
forgotten and $\sigma(v)$ is viewed in the general seminorm spectrum of $R[t]$.
\end{theorem}
\begin{proof}
Both multiplicative functions are the supremum of $|g|_{\sigma(v)}$ for
$g\in I$, giving the first equality. The projective logarithmic function
is $\sup_{g\in I}\log|g|_{\sigma(v)}$. The extended logarithm is an order
isomorphism and preserves all suprema, so this equals
$\log\sup_{g\in I}|g|_{\sigma(v)}$, the intrinsic logarithmic ideal function.
Forgetting the base restriction preserves the underlying seminorm and every
one of its values, so it preserves both suprema and their logarithms.
The boundedness lemma permits the finite generator formulas and ideal arithmetic
of the intrinsic calculus to be applied to these model functions.
\end{proof}

\begin{remark}\label{projective:ideal-bridge-scope}
\uses{projective:ideal-spectrum-bridge}
The correspondence is a formula on polynomial ideals in the normalized model.
Relating it to the source's coherent ideal functions requires the projective
model identification and the homogeneous sheaf data; in the equivariant setting
the appropriate invariance condition is required as well. These geometric
identifications are separate from the intrinsic evaluation equalities proved here.
\end{remark}
